\chapter{Heart Murmur}

\section*{Definition}
An extra or unusual heart sound heard during auscultation, caused by turbulent blood flow across cardiac valves, septal defects, or vascular structures.

\section*{What Happens in Your Body}
Murmurs result from high-flow across normal valves, forward flow across stenotic valves, backward flow across regurgitant valves, or flow across septal defects. Graded by intensity (I-VI), timing (systolic, diastolic, continuous), configuration (crescendo, decrescendo), and quality (blowing, harsh, rumbling).

\section*{Causes}
\begin{itemize}
  \item Innocent (functional) murmur — benign, no structural heart disease
  \item Aortic stenosis (degenerative, bicuspid valve)
  \item Aortic regurgitation
  \item Mitral regurgitation (mitral valve prolapse, papillary muscle dysfunction)
  \item Mitral stenosis (typically rheumatic heart disease)
  \item Pulmonic stenosis or regurgitation
  \item Hypertrophic cardiomyopathy (HOCM)
  \item Ventricular septal defect (congenital)
  \item Atrial septal defect (congenital)
  \item Infective endocarditis with valve vegetation
\end{itemize}

\section*{When to See a Doctor}
See your doctor if:
\begin{itemize}
  \item Symptoms are severe or getting worse over time
  \item New murmur with chest pain, shortness of breath, syncope, fever (possible endocarditis), or unexplained weight loss
  \item You have other concerning symptoms such as fever, weight loss, or bleeding
\end{itemize}

\section*{Related Symptoms}
\begin{itemize}
  \item Shortness of breath
  \item Chest pain
  \item Palpitations
  \item Fatigue
  \item Leg swelling
\end{itemize}
\textbf{Important:} If this symptom persists, your doctor will check for serious underlying causes.

\section*{Self-Care / Home Management}
\begin{itemize}
  \item Rest and avoid strenuous exertion if the murmur is associated with symptoms such as chest pain or shortness of breath.
  \item Maintain good oral hygiene and schedule regular dental visits to reduce the risk of infective endocarditis.
  \item Keep a symptom diary noting palpitations, chest pain, shortness of breath, or leg swelling to share with your doctor.
  \item If you have a known valve abnormality, discuss with your doctor whether antibiotic prophylaxis is needed before dental or surgical procedures.
\end{itemize}

\section*{How Your Doctor Will Evaluate You}
\begin{itemize}
  \item Cardiac auscultation in different positions (supine, sitting, leaning forward, standing) to characterise timing, grade, location, and radiation of the murmur.
  \item Transthoracic echocardiography is the primary imaging tool to assess valve morphology, function, and haemodynamic significance.
  \item Electrocardiogram and chest X-ray are used to evaluate cardiac chamber size, rhythm abnormalities, and pulmonary circulation.
  \item Referral to a cardiologist is standard for newly diagnosed pathological murmurs or for periodic surveillance of known valvular disease.
\end{itemize}

\section*{Treatment Options}
\begin{itemize}
  \item Innocent murmurs require no treatment; reassurance and normal activity are advised with no activity restrictions.
  \item Valve lesions causing symptoms or haemodynamic compromise may require surgical repair or replacement (valvuloplasty, commissurotomy, or prosthetic valve replacement).
  \item Medical management includes diuretics, beta-blockers, or afterload-reducing agents (ACE inhibitors/ARBs) for heart failure symptoms secondary to valvular disease.
  \item Anticoagulation with warfarin is indicated for mechanical prosthetic valves or mitral stenosis with atrial fibrillation; direct oral anticoagulants are not suitable for mechanical valves.
  \item Medications, lifestyle modifications, and physical therapies may be prescribed as appropriate.
  \item Follow-up ensures treatment effectiveness and allows timely adjustments to the care plan.
\end{itemize}

\section*{References}

\begin{thebibliography}{99}
\bibitem{harrison-cardio} Longo DL, et al. \emph{Harrison's Principles of Internal Medicine}. 21st ed. McGraw-Hill; 2022. Chapter 236: Approach to the Patient with a Heart Murmur.
\bibitem{uptodate-murmur} Bonow RO, et al. Auscultation of heart murmurs in adults. In: UpToDate, Post TW (Ed), UpToDate, Waltham, MA; 2025.
\bibitem{otto} Otto CM, et al. \emph{Valvular Heart Disease: A Companion to Braunwald's Heart Disease}. 6th ed. Elsevier; 2023. Chapter 3: Physical Examination.
\bibitem{bailey} Bailey SN, et al. Evaluation of heart murmurs in primary care. \emph{Am Fam Physician}. 2023;107(4):372-381.
\bibitem{bonow2021} Bonow RO, et al. 2020 ACC/AHA guideline for the management of valvular heart disease. \emph{Circulation}. 2021;143(5):e72-e227.
\bibitem{braunwald} Libby P, et al. \emph{Braunwald's Heart Disease}. 12th ed. Elsevier; 2022. Chapter 14: Heart Murmurs and Valvular Disease.
\end{thebibliography}
